\chapter{Experimental Provenance and Historical Configurations}
\label{app:experiment_provenance}

This appendix records the relationship between the archived development
experiments and the final evaluation in Chapter~\ref{chap:experimental_evaluation}.
Historical results remain available for audit, but their runtime scores are
not pooled with or substituted for the final 60 Hz formal results.

\section{Historical preanalysed replay}

The historical runtime suites used file playback with preanalysed musical
timing, a 120 Hz control target, a six-second matching window, a one-second
matching interval and six seconds of warm-up. Runs were headless, with runtime
collision checking enabled. Seeds were 0--4. Durations below include warm-up.
Unlike the final causal protocol, this mode could use musical timing obtained
from audio beyond the current playback position.

\begin{table}[htbp]
\centering\small
\begin{tabular}{p{0.18\linewidth}p{0.56\linewidth}r}
\toprule
Archived suite & Configuration and purpose & Runs \\
\midrule
Replay full & 60 AIST++ music identities and ten CC0 recordings;
full system; 30 s each. Broad playback characterisation. & 350 \\
Replay ablation & 13 inputs, eight conditions, 40 s each.
Exploratory component comparisons within the replay protocol. & 520 \\
Replay long & One repeated stitched input, the same eight conditions,
606 s each. Historical sustained-playback diagnostics. & 40 \\
SMPL subset & 40 music/reference pairs, five seeds, 26 s each.
Traces reused for the separate offline feature analysis. & 200 \\
Preflight & One 10 s development smoke check. & 1 \\
\bottomrule
\end{tabular}
\caption{Archived replay configuration. These 1,111 process runs are separate
from the 135 final formal runtime runs. The SMPL analysis retains its own
200-output denominator.}
\label{tab:app_replay_matrix}
\end{table}

The replay ablation inputs were the stitched AIST++ stream, ten CC0 recordings,
a tempo jump and a silence/recovery stream. Its eight conditions were full,
legacy retrieval, no motion compatibility, instant top-one selection,
no diversity, authored timing, fixed entry/simple transition and no output
limiter. Each condition had $13\times5=65$ runs. Thus a replay ablation
difference used its own matched 65-run full condition, not the separate
350-run Replay full suite. The historical long suite used five runs per
condition, giving $1\times8\times5=40$.

\section{Meaning of historical row names}

\begin{table}[htbp]
\centering\small
\begin{tabular}{p{0.24\linewidth}p{0.67\linewidth}}
\toprule
Former row label & Provenance and configuration \\
\midrule
Replay full, $N=350$ & The archived 70-input playback suite in
Table~\ref{tab:app_replay_matrix}; not a formal 60 Hz result. \\
Legacy retrieval & Replay ablation setting the match policy to legacy. \\
No diversity & Replay ablation setting diversity top-$k$ to 1 and
maximum hold bars to 1,000,000. This is the implemented diagnostic configuration,
not a claim that every possible source of motion diversity is removed. \\
Fixed entry & Replay ablation using the fixed-entry/simple-transition
diagnostic switch; it changes the combined mechanism rather than isolating
entry search alone. \\
Replay authored & Replay ablation using authored motion timing;
its reference is full within that replay ablation suite. \\
Causal authored & A name also used for the earlier 120 Hz causal supplement
in the archived combined analysis. The final 60 Hz condition is explicitly
identified as F-authored in Chapter~\ref{chap:experimental_evaluation}. \\
\bottomrule
\end{tabular}
\caption{Provenance of labels that appeared together in earlier result tables.
Historical replay comparisons and earlier causal comparisons are not arms
of a single formal ablation.}
\label{tab:app_historical_labels}
\end{table}

The earlier causal supplement also had 130 short and five long runs, but ran
at 120 Hz with an earlier implementation. Its 135-run count must not be
confused with the final 135-run cohort. Similarly, the 21-run 120 Hz v5
capacity pilot reported in Chapter~\ref{chap:experimental_evaluation} is a
separate experiment from both the old replay suites and that earlier supplement.
Condition names and equal run counts alone do not identify an experiment.

\section{Permitted comparisons and artifact locations}

Historical replay BAS uses the original preanalysed beat records. Formal
causal BAS uses actual online beat-delivery times. The old records cannot
recover actual frame and response clocks; their nominal-rate joint derivatives
cannot validate the final actual-clock output contract. The old physical-contact
proxy is also distinct from the adapted G1 PFC used in the formal analysis.
Changes in implementation, audio availability, control rate, source coverage
and duration therefore prevent interpreting replay-to-formal differences as
isolated component effects. Reanalysis can improve summaries but cannot recover
measurements that were never recorded.

Within the formal cohort, F-full and F-authored share sources, seeds, durations
and measurement instruments. They support the timing comparison reported in
Chapter~\ref{chap:experimental_evaluation}. F-long supplies sustained-operation
evidence for one repeated source. The historical component ablations are
development diagnostics and do not establish the corresponding component
benefits for the final implementation. Offline retrieval/rejection and SMPL
feature characterisation retain the separate units and limitations stated in
that chapter.

\begin{samepage}
For reproducibility, the following paths are relative to the repository root:
\begin{itemize}[leftmargin=*]
    \item Archived replay commands, manifests and per-suite run indices:
    \path{realtime/humanoid_robot/src/test/output/thesis_final/}.
    \item Earlier causal supplement:\\
    \path{realtime/humanoid_robot/src/test/output/thesis_supplement/}.
    \item Archived combined replay/causal analysis containing the former
    mixed table rows: \path{docs/thesis/experiment_results/reanalysis_v2/}.
    \item Separate 120 Hz v5 capacity-pilot analysis:
    \path{docs/thesis/experiment_results/reanalysis_pilot_120_v5/}.
    \item Final formal runtime commands, manifest and run indices:
    \path{realtime/humanoid_robot/src/test/output/thesis_rescue_formal_60_v5/}.
    \item Final analysis and the tables/macros imported by the thesis:
    \path{docs/thesis/experiment_results/reanalysis_rescue_60_v5/}.
\end{itemize}
\end{samepage}
The final analysis bundle retains historical material for provenance and
offline feature reconstruction. The runtime table export explicitly selects
its final \texttt{online\_causal} short and long groups. Presence in that bundle
does not make an archived replay run part of the formal cohort. Each cohort's
manifest and recorded commands determine its implementation and configuration.
